% =====================================================================
%  31.2  Prim's algorithm and the frontier view.
%  Exports: sec:c17-prim  thm:c17-prim  lem:c17-frontier
%  def:c17-frontier  lst:c17-prim  lst:c17-prim-matrix  fig:c17-prim-trace
%  WORKED: a full Prim run on the running graph G, recovering the same MST
%  {1-2, 3-4, 1-3, 4-5} of weight 12 that Kruskal found in 31.1, so the two
%  traces agree.  No new solved problem: this section is the dense-graph twin
%  of 31.1 and the "Prim is not Dijkstra" disambiguation.
% =====================================================================
\section{Prim's Algorithm and the Frontier View}\label{sec:c17-prim}

\begin{previouslybox}[Ch.31.1 grew the MST edge by edge, globally sorted]
\(\secref{c17-kruskal}\) sorted every edge once and unioned components.  That is
ideal when \(m\) is small.  When the graph is dense -- or given as a weight
matrix -- sorting \(\Theta(n^2)\) edges is wasteful.  Prim reaches the same tree
by growing a \emph{single} component and always taking the cheapest edge leaving
it, touching only the frontier.
\end{previouslybox}

\begin{topicintro}
\term{Prim}'s algorithm maintains one growing tree and a \term{frontier}: the set
of edges with exactly one endpoint inside the tree.  At each step it takes the
lightest frontier edge (safe by the cut property, since the tree-vertices versus
the rest is a cut and that edge is its cheapest crossing).  A priority queue keyed
by edge weight gives \(\bigO{m\log n}\); on a dense graph an array-based
\(\bigO{n^2}\) scan wins instead.

The section's other job is a disambiguation.  Prim and Dijkstra look almost
identical -- both pop a vertex and relax to a frontier priority queue -- but they
relax on \emph{different} quantities.  Confusing them is the single most common
Prim bug, so we state the distinction sharply and then run Prim by hand on the
running graph \(G\).

\textbf{Prerequisites.} \(\secref{c17-kruskal}\) -- the cut property, which is
Prim's correctness proof verbatim.  A priority queue
(\texttt{std::priority\_queue}).  Nothing
about DSU is needed here; Prim never merges components, it only grows one.

\textbf{After this section you can}
\begin{itemize}[nosep]
  \item run Prim with a heap in \(\bigO{m\log n}\) and with a matrix in \(\bigO{n^2}\); \hfill\lv{4}
  \item explain precisely why Prim's relaxation differs from Dijkstra's; \hfill\lv{3}
  \item choose Kruskal vs Prim from the graph's density. \hfill\lv{4}
\end{itemize}
\end{topicintro}

% ---------------------------------------------------------------
\subsection{Growing one component from a frontier}
% ---------------------------------------------------------------

\begin{definition}{Frontier (cut) of a partial tree}{c17-frontier}
Let \(S\subsetneq V\) be the vertices already in the growing tree.  The
\term{frontier} is the set of edges \(\{(u,v): u\in S,\ v\notin S\}\); it is
exactly the set of edges crossing the cut \((S, V\setminus S)\).
\end{definition}

\begin{theorem}{Prim's algorithm}{c17-prim}
Start from any vertex.  Repeatedly add the minimum-weight frontier edge and its
outside endpoint.  When all vertices are in the tree, the chosen edges form an
MST.
\end{theorem}

\begin{proofbox}
At each step the frontier is the crossing set of the cut \((S,V\setminus S)\), so
by the cut property (\(\thmref{c17-cut}\)) its cheapest edge lies in some MST;
adding it is safe.  After \(n-1\) steps we hold \(n\) vertices and \(n-1\) safe
edges forming a tree, hence an MST. \(\Box\)
\end{proofbox}

\begin{lemma}{Prim is not Dijkstra}{c17-dijkstra}
Both maintain a frontier priority queue, but Prim stores for each outside vertex
the \emph{single cheapest edge} from the tree to it (key \(=\min_{u\in S} w(u,v)\)),
while Dijkstra stores the \emph{accumulated path length} from the source
(key \(=\mathrm{dist}(s,u)+w(u,v)\)).  Prim's key never adds the cost to reach
\(u\); it only ever measures one edge.
\end{lemma}

\begin{lstlisting}[caption={Prim with a binary heap. Push (weight, to) neighbours of each newly added vertex; lazily skip stale entries whose endpoint is already in the tree. Costs O(m log n). Returns MST weight and edge count}, label={lst:c17-prim}]
#include <bits/stdc++.h>
using namespace std;
typedef long long ll;
typedef pair<ll,int> pli;
struct Adj { int to; ll w; };
vector<vector<Adj> > g;
vector<char> inTree;
ll prim(int n){
  inTree.assign(n+1, 0);
  priority_queue<pli, vector<pli>, greater<pli> > pq;   // min-heap (w, vertex)
  pq.push({0, 1});                                       // start from vertex 1
  ll total=0; int used=0;
  while(!pq.empty() && used<n){
    ll w; int v; tie(w,v)=pq.top(); pq.pop();
    if(inTree[v]) continue;                             // stale frontier entry
    inTree[v]=1; total+=w; used++;
    for(const Adj& e : g[v]) if(!inTree[e.to])           // widen the frontier
      pq.push({e.w, e.to});
  }
  return (used==n) ? total : -1;
}
\end{lstlisting}

\begin{warningbox}[The Dijkstra reflex]
The temptation is to push \(\{w + \mathrm{accumulated},\ e.to\}\) into the heap.
That is Dijkstra and it is \emph{wrong} for the MST: Prim's frontier key is the
edge weight alone.  A clean discriminator: on the triangle \(s\!-\!a{:}2,\ s\!-\!b{:}3,\ a\!-\!b{:}1\),
Prim starts at \(s\), takes \(s\!-\!a{:}2\), then sees frontier \(\{s\!-\!b{:}3,\ a\!-\!b{:}1\}\)
and takes the cheap \(a\!-\!b{:}1\) -- MST \(\{s\!-\!a,\ a\!-\!b\}\), weight \(3\).
Dijkstra instead routes \(b\) by accumulated distance and would hang \(b\) straight
off \(s\) (path \(3\) beats \(2+1=3\) only on a tie it resolves toward the source
edge), yielding the shortest-path tree \(\{s\!-\!a,\ s\!-\!b\}\), weight \(5\).  Same
frontier machinery, different key -- never borrow the other's relaxation.
\end{warningbox}

% ---------------------------------------------------------------
\subsection{The dense-graph \texorpdfstring{\(\bigO{n^2}\)}{O(n^2)} variant}
% ---------------------------------------------------------------

When the graph is given as a matrix (or \(m=\Theta(n^2)\)), the heap spends most
of its time churning dense frontier updates.  A plain scan for the cheapest
outside vertex each round is \(\bigO{n^2}\) total and beats the heap for
\(n\lesssim 5000\): keep an array \(\mathit{best}[v]\) of the cheapest edge from
the tree to each outside vertex, pick the minimum \(\mathit{best}\), then relax
\(\mathit{best}\) with the new vertex's row.

\begin{lstlisting}[caption={Matrix Prim, O(n squared). The best array holds the cheapest tree-to-vertex edge weight; pick the minimum over outside vertices, then relax it with the newly added row. Ideal for dense graphs or an explicit weight matrix}, label={lst:c17-prim-matrix}]
#include <bits/stdc++.h>
using namespace std;
typedef long long ll;
ll w[5005][5005];
ll primMatrix(int n){
  vector<ll> best(n+1, LLONG_MAX); vector<char> done(n+1,0);
  best[1]=0; ll total=0;
  for(int it=1; it<=n; it++){
    int v=-1;
    for(int i=1;i<=n;i++) if(!done[i] && (v==-1 || best[i]<best[v])) v=i;
    if(best[v]==LLONG_MAX) return -1;             // disconnected
    done[v]=1; total+=best[v];
    for(int to=1; to<=n; to++)
      if(!done[to] && w[v][to] < best[to]) best[to]=w[v][to];  // relax frontier
  }
  return total;
}
\end{lstlisting}

% ---------------------------------------------------------------
\subsection{WORKED: Prim on the running graph}
% ---------------------------------------------------------------

\begin{tracebox}[Prim (heap) on \(G\), starting at vertex 1]
Frontier keyed by the single cheapest tree-to-outside edge.
\begin{itemize}[nosep,leftmargin=1.4em]
  \item Tree \(=\{1\}\). Frontier: \(1\!-\!2{:}1,\ 1\!-\!3{:}3\). Cheapest
        \(1\!-\!2{:}1\) -- take, add 2.
  \item Tree \(=\{1,2\}\). Frontier: \(1\!-\!3{:}3,\ 2\!-\!3{:}4,\ 2\!-\!4{:}5\).
        Cheapest \(1\!-\!3{:}3\) -- take, add 3. (Note \(2\!-\!3{:}4\) is a
        pricier edge to the same vertex 3; only the cheapest is keyed.)
  \item Tree \(=\{1,2,3\}\). Frontier: \(3\!-\!4{:}2,\ 2\!-\!4{:}5,\ 3\!-\!5{:}7\).
        Cheapest \(3\!-\!4{:}2\) -- take, add 4.
  \item Tree \(=\{1,2,3,4\}\). Frontier: \(2\!-\!4{:}5\) (dead, 4 in tree),
        \(4\!-\!5{:}6,\ 3\!-\!5{:}7\). Cheapest \(4\!-\!5{:}6\) -- take, add 5.
  \item All five vertices in the tree.
\end{itemize}
Chosen edges \(\{1\!-\!2,\ 1\!-\!3,\ 3\!-\!4,\ 4\!-\!5\}\), weight
\(1+3+2+6=\mathbf{12}\) -- the \emph{same} MST as \(\secref{c17-kruskal}\), reached
in a different order (Prim picked \(1\!-\!3{:}3\) before \(3\!-\!4{:}2\); Kruskal
took them sorted).
\end{tracebox}

\begin{figure}[H]
\centering
\begin{tikzpicture}[diagram, scale=0.9, transform shape,
   tnode/.style={circle, draw=primaryblue!70, fill=primaryblue!8, line width=0.7pt,
                 minimum size=7mm, font=\scriptsize\bfseries, inner sep=1pt},
   acc/.style={line width=1.6pt, draw=successgreen!80!black},
   ord/.style={circle, fill=deepblue, text=white, inner sep=1.2pt, font=\scriptsize\bfseries}]
  \node[tnode] (v1) at (0,0) {1};
  \node[tnode] (v2) at (3,1.2) {2};
  \node[tnode] (v3) at (3,-1.4) {3};
  \node[tnode] (v4) at (6,-0.4) {4};
  \node[tnode] (v5) at (8.2,1) {5};
  \draw[acc] (v1)--node[ord] {1} (v2);
  \draw[acc] (v1)--node[fill=white,inner sep=1pt,font=\scriptsize,pos=0.7] {3} (v3);
  \draw[acc] (v3)--node[ord] {2} (v4);
  \draw[acc] (v4)--node[ord] {3} (v5);
  \node[annot, font=\scriptsize] at (4,-2.6) {badges = Prim take order; same four green edges as Kruskal};
\end{tikzpicture}
\caption{Prim's take order on \(G\).  The badge on each edge is the step (1--4)
that added it; edge \(1\!-\!3\) carries its weight \(3\) instead since it shares
the round with vertex 3's insertion.  The final tree equals Kruskal's.}
\label{fig:c17-prim-trace}
\end{figure}

\begin{comparisonbox}[Kruskal vs Prim: which to run]
\begin{center}\renewcommand{\arraystretch}{1.2}\setlength{\tabcolsep}{4pt}
\begin{tabular}{@{}l l l@{}}
\toprule
 & \textbf{Kruskal} & \textbf{Prim (heap)}\\
\midrule
Time & \(\bigO{m\log m}\) & \(\bigO{m\log n}\)\\
Dense \(m=\Theta(n^2)\) & sort-heavy & \(\bigO{n^2}\) array wins\\
Sparse & default choice & fine but no edge\\
Needs & a DSU (\chref{c28}) & a priority queue\\
Grows & many components & one component\\
\bottomrule
\end{tabular}\end{center}
Use Kruskal when edges are already a list and the graph is sparse; use Prim
(matrix form) when the graph is dense or arrives as an adjacency matrix.  Both
return the identical MST weight; pick for the constant factors.
\end{comparisonbox}

\begin{recallbox}[Recall -- what \(\secref{c17-prim}\) gave you]
\selftest{
  \item What cut does Prim's cheapest-frontier-edge appeal to for correctness?
  \item State the one-line difference between a Prim key and a Dijkstra key.
  \item For \(n\le 5000\) dense graphs, which Prim implementation wins and why?
  \item Did Prim on \(G\) pick the same four edges as Kruskal? In what order?
}
\end{recallbox}
